\section{Conclusiones e integración con E4}
\label{sec:conclusiones}

Este trabajo abordó dos dimensiones de la Entrega 3 sobre la base de código de microservicios
construida en entregas anteriores: la fragmentación horizontal correcta y verificable de la capa
de datos, y el procesamiento analítico paralelo de un volumen de datos que excede lo razonable
para un solo proceso. Sobre la primera dimensión, se rediseñó la tabla \texttt{tickets} de un
esquema fragmentado por zona geográfica a uno fragmentado por rango de \texttt{fecha\_apertura},
verificando formalmente las tres condiciones de corrección de toda fragmentación horizontal
—completitud, reconstrucción y disyunción— y documentando explícitamente, en vez de ocultar, los
\emph{trade-offs} aceptados: degradación de las consultas por zona a \emph{scatter-gather} sobre
cuatro particiones, y la imposibilidad de colocalización física con la tabla de clientes por
pertenecer a un microservicio distinto. La corrección del comportamiento transaccional del clúster
—en particular, el rechazo explícito de escrituras concurrentes conflictivas propio del aislamiento
\texttt{SERIALIZABLE}— se verificó empíricamente contra una instancia real de CockroachDB mediante
Testcontainers, no mediante inspección manual ni simulación, y se instrumentó con métricas
Prometheus validadas sin advertencias por el \emph{linter} oficial de la herramienta.

Sobre la segunda dimensión, se implementó un pipeline de Apache Spark con las cinco categorías de
transformación exigidas, aplicadas de forma coherente a un objetivo analítico propio del dominio
—reincidencia de clientes y agrupamiento de incidencias por texto y zona— y ejecutado de extremo a
extremo sobre un conjunto sintético de 600{,}000 registros con resultados reproducibles y no
triviales: 68.3\% de reincidencia detectada dentro de una ventana de 30 días y una distribución de
\emph{clusters} no degenerada. Siguiendo metodología estándar de experimentación en ingeniería de
software, se ejecutó además un protocolo formal de 50 corridas contrastando este pipeline contra un
baseline secuencial en pandas, con un resultado que no era el esperado y que precisamente por eso
es valioso: para este volumen de datos, el baseline secuencial resultó significativamente más
rápido que Spark incluso en su configuración de mayor paralelismo
($p \approx 5\times10^{-10}$), aunque Spark sí exhibe un \emph{speedup} interno real al escalar de
uno a ocho núcleos (ajuste de Amdahl $p \approx 0.53$). Reportar este resultado tal como ocurrió,
en vez de solo buscar el resultado que confirmara la hipótesis inicial, es coherente con el
principio de honestidad aplicado en todo este documento.

La caracterización de tolerancia a fallos ante la caída de uno y dos nodos, diseñada e
implementada en la Sección~\ref{sec:cluster-verificacion}, se ejecutó y midió por completo: el
clúster tolera la caída de un nodo sin pérdida de datos ni de disponibilidad de escritura (0
errores en 691 filas, pico de latencia P95 de 1.59\,s durante la reelección de líder), y pierde
disponibilidad de escritura durante 62 segundos al caer dos de tres nodos, recuperándose de
inmediato al reincorporar un nodo — ambos resultados consistentes con el quórum de mayoría
exigido por el algoritmo de consenso Raft descrito en la Sección~\ref{sec:fundamento-teorico}.

\subsection{Integración con la Entrega 4}

Siguiendo el principio de complementariedad E3→E4 de la guía de esta entrega —todo artefacto
producido en E3 debe reutilizarse, ampliarse o instrumentarse en E4, no descartarse—, los
siguientes artefactos quedan disponibles para consumo directo en el siguiente ciclo:

\begin{itemize}[leftmargin=1.5em]
    \item \textbf{El esquema de base de datos fragmentado} (\texttt{db-cluster/scripts/init\_db.sql},
        ADR-0003) y el propio \textbf{clúster de tres nodos de CockroachDB} quedan como la capa de
        persistencia definitiva del sistema: E4 no necesita rediseñar ni migrar datos, solo construir
        sobre ella la interfaz web, la interfaz móvil y la observabilidad exigidas en el siguiente
        ciclo.
    \item \textbf{Las tres métricas Prometheus} instrumentadas en la
        Sección~\ref{sec:cluster-verificacion} (\texttt{crdb\_query\_duration\_seconds},
        \texttt{crdb\_transaction\_retries\_total}, \texttt{crdb\_pool\_active\_connections}) se
        emplean tal cual, sin modificación, como fuente de datos del panel de observabilidad
        consolidado que exige E4 —por eso se instrumentaron en esta entrega y no se aplazaron.
    \item \textbf{Las pruebas de integración con Testcontainers} quedan incorporadas a la suite de
        pruebas del proyecto como regresión permanente: cualquier cambio futuro sobre
        \texttt{ticket-service} en E4 se sigue verificando contra una instancia real de CockroachDB,
        no contra una base embebida.
    \item \textbf{El pipeline de Spark y el protocolo experimental} (\texttt{spark/src/pipeline.py},
        \texttt{baseline.py}, \texttt{protocol\_experiment.py}) quedan como la base analítica
        reutilizable: E4 puede extenderlos con nuevas transformaciones o exponer sus resultados
        (reincidencia, \emph{clusters} de incidencias) a través de la interfaz de usuario, sin
        rehacer la integración JDBC ni el diseño experimental ya validado aquí.
\end{itemize}

Como trabajo futuro más allá de lo directamente heredable por E4, se identifican dos líneas
adicionales: una estrategia de fragmentación compuesta que preserve el beneficio de poda de
particiones tanto para consultas temporales como para consultas por zona, y la validación del
pipeline analítico sobre un conjunto de datos real de incidencias de un ISP, para contrastar la
representatividad del dataset sintético usado en esta entrega frente a la distribución real de
texto y de cardinalidad de clientes del dominio.
